\documentclass[
oneside,							% einseitiges Layout
captions=signature,	  % caps immer unter Element (Bild, Grafik...)
openany,							% Kapitel können auch auf linken Seiten beginnen
listof=totoc				% Abb.- und Tab.verzeichnis im Inhalt	
]
{scrbook}


\usepackage[utf8]{inputenc}
\usepackage[T1]{fontenc}

\usepackage{mathptmx}
\usepackage[scaled=.90]{helvet}
\usepackage{courier}
					
\usepackage[ngerman]{babel}

\usepackage{tikz}
\usepackage{qtree} % Baumstrukturen zeichnen

\usepackage{tabularx}
\usepackage{tabulary}
\usepackage{multirow}
\usepackage{colortbl}
\usepackage{booktabs}
\usepackage{listings}

%% Packages für Grafiken & Abbildungen %%%%%%%%%%%%%%%%%%%%%%

\usepackage{graphicx} 	% Zum Laden von Grafiken
\usepackage{xcolor}
\usepackage{color}
\usepackage{microtype} 	% visual enhancement randausgleich
\usepackage{float} 			% ermöglicht den H param bei Bilder 
												% für Plazierungen an der Position

\usepackage{rotating} % Text drehen

\usepackage[nogroupskip, nonumberlist, toc]{glossaries}
\usepackage{glossary-inline}
\usepackage[printonlyused]{acronym}

\usepackage{pifont} 	%symbol fonts like Windings
\usepackage{amsmath}	% Math font optimisation
\usepackage{amssymb} 	% Math symbols

\usepackage{tocloft} % Table of content - list of figures, tables - management


%% Bibliographiestil %%%%%%%%%%%%%%%%%%%%%%%%%%%%%%%%%%%%%%%%%%%%%%%%%%

\usepackage{natbib}
\renewcommand{\familydefault}{\sfdefault}


\setlength{\parindent}{0cm}			% Absatzeinzug
\setlength{\parskip}{6pt}				% Leerzeile zwischen Absätzen im Inhaltsverzeichnis
\textwidth=15.7cm								% Textbreite
\leftmargin=-3cm

\textheight=22.6cm							% Texthöhe
\topmargin=-0.7cm								% Rand (Vorsicht minus)
\headheight=0.2cm								% Höhe der Kopfzeile
\headsep=1.5cm									% Abstand von der Kopfzeile zum Text

%\setlength{\bibsep}{0.4\parskip}   	 	%Bibliographieabstand fest
\setlength{\bibsep}{0.3em plus 0.5em}	
%Bibliographieabstand variabel um Seitenumbruch zu unterdrücken

\setlength{\bibhang}{1cm}    			%Bibliographieeinzug
%\newcommand{\bibfont}{\normalsize}   	%Bibliographiezeichengröße
\renewenvironment{quote}{\list{}{\rightmargin0pt\leftmargin1cm}\item\relax}{\endlist} 
% neue \quote-Definition als globale Variabel für rechtsbündige Zitate

\setcitestyle{notesep={: }}
\setcitestyle{aysep={}}

%%%%%%%%%%%%%%%%%%%%%%%%%%%%%%%\setlist{\setlength{\itemsep}{0.8em plus 2em}, parsep=.65em, partopsep=1em} %listenformatierung

%%%%%% TOC Tiefe %%%%%%%%%%%%%%%%%%%%%%%%%%%%%%%%%%%%%
\setcounter{tocdepth}{3}

%%%%%%%% Nummerierungstiefe absätze %%%%%%%%%%%%%%%%%%
\setcounter{secnumdepth}{3}

%\setkomafont{section}{\LARGE}
%\setkomafont{subsection}{\Large}
%\setkomafont{subsubsection}{\large}
%\setkomafont{paragraph}{\large}
%\setkomafont{subparagraph}{\large}
% Die Überschriften sind in der Lokalen *.cls 
% Classendefinition definiert %%%%%


%% Abk. und Kapitälchen%%%%%%%%%%%%%%%%%%%%%%%%%%%%%%%%%%%%%%%%%%%%%%%%


\setglossarystyle{index}
\makenoidxglossaries
\setacronymstyle{sc-short-long}

\newacronym{fmx}{Formex}{Formalized Exchange of Electronic Publications}
\newacronym{cl}{cl}{Computerlinguistik der Universität Trier}
\newacronym{darpa}{darpa}{Defence Advanced Research Projects Agency}
\newacronym{bleu}{bleu}{Bilingual Evaluation Understudy}
\newacronym{nist}{nist}{National Institute of Standards and Technology}
\newacronym{hter}{hter}{Human-targeted Translation Edit Rate}
\newacronym{rbmt}{rbmt}{Rule Based Machine Translation -- Regelbasierte Maschinelle Übersetzung}
\newacronym{smt}{smt}{Statistical Machine Translation -- Statistische Maschinelle Übersetzung}
\newacronym{meteor}{meteor}{Metric for Evaluation of Translation with Explicit ORdering}
\newacronym{lepor}{lepor}{Length Penalty, Precision, n-gram Position difference Penalty and Recall}
\newacronym{amber}{amber}{A Modified Bleu Enhanced Ranking}

\newglossaryentry{glos:precision}{name=precision, description={Genauigkeit}}
\newglossaryentry{glos:recall}{name=precision, description={Trefferquote}}


%%% Trennungsregeln für das gesamte Dokument.

\hyphenation{Über-set-zungs-epo-che}
\hyphenation{Über-set-zungs-ver-fah-ren}
\hyphenation{Über-set-zung}
\hyphenation{Kor-re-la-tions-koef-fi-zient}
\hyphenation{Kor-re-la-tions-koef-fi-zien-ten}
\hyphenation{Rang-Kor-re-la-tions-koef-fi-zien-ten}


%%definition von quotation

\let\oldquotation\quotation
\renewcommand\quotation{\oldquotation\noindent\renewcommand{\baselinestretch}{1}\small\ignorespaces} 

%% Absolute Value definition %

\delimitershortfall-1sp
\newcommand\abs[1]{\left|#1\right|}
%%%%%%%%%%%%%%%%%%%%%%%%%%%%%%%%%%%%%%%%%%%%%%%

\providecommand\phantomsection{}

% Farbige unterteilung

\newcommand{\status}[2]{%
  \colorlet{author#1}{#2}% Create an author colour
  \expandafter\def\csname authoredby#1\endcsname{% Create author colour settings
    \renewcommand{\cftchapfont}{\bfseries\color{author#1}}% Chapter colour
    \renewcommand{\cftsecfont}{\color{author#1}}% Section colour
    \renewcommand{\cftsubsecfont}{\color{author#1}}% Subsection colour
	\renewcommand{\cftsubsubsecfont}{\color{author#1}}}% Subsection colour
	}

\makeatletter
\newcommand{\setstatus}[1]{\addtocontents{toc}{\protect\@nameuse{authoredby#1}}}%
\makeatother

\status{R}{red}% Author A will be coloured red
\status{Y}{orange}% Author B will be coloured blue
\status{G}{green}% Author B will be coloured blue


\begin{document}




%%%%%% T I T E L S E I T E %%%%%%%

\thispagestyle{empty}

\vspace*{-2cm} \hspace{-1cm}
{\includegraphics[height=2.5cm]{BILDER/logo.png}} \\ \vspace{3cm}
	\begin{center}
		\textsf{ \Huge{\textbf{Hausarbeit}}}\\ \vspace{1.5cm}
		\textsf{ \Huge{zu dem Thema}} \vspace{1.5cm}
			\begin{center}
				\textsf{ \Huge{\textbf{Vorlage zur Erstellung einer Hausarbeit im Fach CL}}}
			\end{center}
			\vfill \vfill \textsf{\Large \today}
	\end{center}
\vfill \hspace{-2cm}
\textsf{{\large
\begin{tabular}{ll}
		Vorname Name\\
		Email: mein.name@uni-trier.de\\
\end{tabular}
}}
\newpage




%% Erzeugung von Verzeichnissen %%%%%%%%%%%%%%%%%%%%%%%%%%%%%%%%%%%%%%%
\addtocontents{toc}{\protect\thispagestyle{empty}} % unterdrückt Seitenzahl
\addtocontents{lot}{\protect\thispagestyle{empty}} % unterdrückt Seitenzahl
\addtocontents{lof}{\protect\thispagestyle{empty}} % unterdrückt Seitenzahl

\tableofcontents			% Inhaltsverzeichnis



%% Der Text %%%%%%%%%%%%%%%%%%%%%%%%%%%%%%%%%%%%%%%%%%%%%%%%%%%%%%%%%%%
\frontmatter					% Vorspann (z.B. Glossar)

\phantomsection 

\printnoidxglossary[sort=letter, title=Abkürzungsverzeichnis, toctitle=Abkürzungsverzeichnis]

\mainmatter						% Hauptteil

\renewcommand{\baselinestretch}{1} %%% Zeilenabstand für das Dokument definieren.



\chapter{Einleitung}
\label{sec:Einleitung}

"`Augusta"' ist ein Chatbot, dessen Aufgabe es ist, dem Kunden Produkte aus dem Uni-Shop zu empfehlen. Dafür muss der Nutzer dem Chatbot via Texteingabe angeben, was gesucht wird. Diese Nutzereingaben werden mit Hilfe von Pattern-Matching analysiert und via PostgreSQL-Schnittstelle werden Antworten basierend auf Queries generiert, die dem geäußerten Wunsch des Nutzers entsprechen.\\
Das Ergebnis des Projekts kann auf GitHub unter \url{https://github.com/AlKYala/Chatbot} eingesehen werden.\\


\section{Ziele des Projekts}
\label{sec:ZielDesProjekts}

Vorrangiges Ziel des Projekts war der Versuch der Entwicklung eines computerlinguistischen Programms. Das Umsetzen von im Studium angeeigneter Fähigkeiten theoretischer und praktischer Natur, wie beispielsweise der Automatentheorie, Datenbankentechnologie und Softwaretechnik, hat eine essentielle Rolle in der Konzeption und Realisierung des Programms gespielt. Weiteres Ziel war das Einarbeiten in ein fremdes Framework mit eigenem Interpreter, hier ChatScript \citep{chatscript2019} von Bruce Wilcox, um ein entsprechendes Programm zu entwickeln.\\


\section{Der Chatbot Augusta}
\label{sec:Augusta}

Im Projektseminar der Computerlinguistik haben wir in kleinen Gruppen jeweils einen Chatbot für eine spezifische Anwendung entwickeln wollen. Unsere Gruppe, bestehend aus Ali Yalama und Julia Roegner, hat sich dazu entschieden, einen Chatbot als Kaufberatung für den Uni-Shop der Universität Trier zu bauen. Der Bot sollte für das Deutsche gelten und eine ausschließlich beratende Funktion haben, das bedeutet, dass der Bot Bestellungen oder Fragen zum Uni-Shop nicht beantworten können muss.\\
Als Basis dienten die Beispieldialoge (s. Anhang), die auch am Ende dieser Dokumentation abgedruckt sind. Hier wird davon ausgegangen, dass der Nutzer noch nicht weiß, welches Produkt er haben möchte, sondern ihm nur bekannt ist, wofür er es möchte, also entweder eine Geschenkidee hat oder etwas für eine bestimmte Anwendung sucht. Dann bietet der Chatbot dem Nutzer die Kategorien an, die in der Datenbank für Anwendungsgebiet oder Geschenkidee hinterlegt sind. Im Weiteren kann der Kunde dem Chatbot noch weitere Informationen über den gesuchten Artikel wie eine Farbe oder den Preis übergeben, am besten in Form eines kompakten Satzes, den der Chatbot analysiert. Dann nutzt der Bot die extrahierten Stichwörter, um in einer Datenbank, in der die Informationen über die Produkte des Uni-Shops gespeichert sind, nach etwas Passendem zu suchen.\\
Der gefundene Gegenstand wird dem Kunden angeboten und der Kunde kann entscheiden, ob er ihn ablehnt, dann sucht der Bot erneut, oder ihn annimmt, woraufhin der Chatbot fragt, ob eine neue Suche mit anderen Kriterien durchführen soll. Wenn nicht, beendet sich der Bot.\\
Der Name des Bots, Augusta, leitet sich vom Namen der antiken Stadt \textit{Augusta Treveorum} an der Mosel her, aus der das heutige Trier hervorgegangen ist.\\


\section{Struktur der Dokumentation}
\label{sec:Inhaltsbeschreibung} 
In Kapitel \ref{sec:InstallationAugusta} werden die einzelnen Versionen Augustas erläutert sowie die Installation beschrieben. In Kapitel \ref{sec:ChatScript: Grundlagen} folgen dann einige grundlegende Regeln zum Schreiben von Chatbots in ChatScript, woran sich ein Kapitel zu der Programmlogik hinter unserem Chatbot anschließt, das Augusta als deterministischen Automaten vorstellt. Informationen zu der Datenbank \textit{Uni-Shop}, in der die Produktinformationen der Artikel gespeichert sind, und ihrer Einbindung sowie der Suchanfragen finden sich in Kapitel \ref{sec:DBuEinbindung}. Auf die Benutzerschnittstelle wird in Kapitel \ref{sec:Benutzerschnittstelle} eingegangen. In Kapitel \ref{sec:Probleme} werden schließlich einige Probleme erläutert, auf die wir gestoßen sind und im Rahmen des Projekts leider nicht mehr lösen konnten.\\
Am Ende der Dokumentation sind die Dialoge abgedruckt, die als Basis dienten, und die Zusammenfassung des Ergebnisses aus einem kleinen Usability Test.

\chapter{ChatScript: Grundlagen}
\label{sec:ChatScript: Grundlagen}

Bei ChatScript handelt es sich um ein Framework, der es Entwicklern erlaubt, regelbasierte Chatbots zu entwickeln. Es handelt sich hierbei um "Natural Language tool" mit eigenem Interpreter, welcher Code in C++ übersezt. 


 

\section{ChatScript: Regeln}
\label{sec:ChatScript: Regeln}

\section{ChatScript: Topics}
\label{sec:ChatScript: Topics}

Die Quelldateien in einem ChatScript-Programm werden auch als "Topics" bezeichnet und haben die Dateiendung .top. Topics können als Zusammenfassung mehrerer Regeln betrachtet werden. Es bietet sich an, Topics als Module eines Chatbots zu handhaben. So wird z.B. in diesem Projekt "Augusta" die Datenbankabfrage als eignenes Topic gehandhabt. 

\section{ChatScript: Konzepte}
\label{sec:ChatScript: Konzepte}

\section{ChatScript: Bedingungen zur Steuerung von Programmablauf}
\label{sec:ChatScript: Bedingungen zur Steuerung von Programmablauf}

\section{ChatScript: Befehle zur Steuerung von Programmablauf}
\label{sec:ChatScript: Befehle zur Steuerung von Programmablauf}

\section{ChatScript: Pattern-Matching}
\label{sec:ChatScript: Pattern-Matching}

\section{ChatScript: Zufällige Ausgabe}
\label{sec:ChatScript: Zufällige Ausgabe}

\chapter{Augusta: Programmlogik}
\label{sec:Augusta: Programmlogik}



Hier, stelle in z.B. einem Automaten dar, wie der Programmablauf normalerweise ablaufen soll
\section{Befehle und Programmablauf}
\label{sec:Befehle und Programmablauf}

\section{Variablen und Programmablauf}
\label{sec:Variablen und Programmablauf}

\section{Konzepte und Programmablauf}
\label{sec:Konzepte und Programmablauf}

\section{Programmablauf im Detail}
\label{sec:Programmablauf im Detail}


\chapter{Kapitel1j}
\label{sec:Kapitel1j}

\section{UnterKapitel2j}
\label{sec:UnterKapitel2j}

\section{UnterKapitel3j}
\label{sec:UnterKapitel3j}
\chapter{Probleme}
\label{sec:Probleme}

Dieses Kapitel behandelt Schwierigkeiten, die bei der Arbeit mit ChatScript aufgetaucht sind und haben für uns ein ernstes und großes Hindernis dargestellt. Im Falle des zweiten Problems haben wir auch keine Lösung gefunden.\\


\section{Zwei-Query-Problem}
\label{Section: Zwei-Query Problem}

Ein Problem, für das ein Workaround gefunden werden konnte, war das 'Zwei-Query-Problem'. Dieses Problem ist in den allen Versionen vorhanden und tritt wegen einem der folgenden Fehler auf:

\begin{enumerate}
\item {Suche auf einer vor der Nutzung existierenden Datenbank}
\item {Fehlschlag in Queries}
\end{enumerate}  

Ersteres tritt auf, wenn man mit Hilfe des Befehls \lstinline|:build onesentence| oder Ähnliches den Bot, hier onesentence, neu kompiliert. Dies hat den Ursprung darin, dass der Bezeichner für die Datenbanken per Client mit einer Zufallszahl konkateniert wird und jeder \lstinline|:build|-Befehl dieselbe Zufallszahl generiert. Wird eine Tabelle nach ihrer Nutzung nicht gelöscht, dann tritt dieser Fehlschlag auf.\\
Dieses Problem erscheint in dieser Form nicht außerhalb des Testens, da es unwahrscheinlich ist, dass zwei Tabellen mit demselben Bezeichner zufallsgeneriert werden. Die Komplikation entsteht nicht, wenn ein Bot regelmäßig, das heißt durch Verabschiedung in Ende, beendet wird. Man kann dieses Problem manuell durch explizite \lstinline|drop|-Befehle außerhalb von Chatscript behandeln.\\
Zweiteres tratt auf, wenn mit den angegebenen Eigenschaften kein Produkt gefunden werden konnte. Wenn kein Produkt gefunden werden kann, ist dem Skript zu Folge in den Zustand \texttildelow ende.SETNULLANDREP überzugehen, um eine neue Suche zu gewährleisten. Nach der Programmlogik müssten hier die einzelnen Regeln bis zu der Datenbanksuche wieder iteriert werden, um Nutzerangaben zu sammeln. Stattdessen versuchte der Bot nach dem Zurücksetzen der Variablen, erneut eine Suche mit NULL-Variablen auszuführen, indem er eigenständig in SUMMARY sprang und ohne eine Antwort abzuwarten eine Suche startete.\\
Dieses Problem wurde größtenteils umgangen, indem man den Nutzer mehrmals hintereinander fragt, ob man eine neue Suche ausführen will. Der Ursprung dieses Problems bleibt jedoch ungeklärt. Eine Vermutung ist, dass ChatScript intern Regeln, die bereits besucht worden sind, markiert, um nicht erneut von alleine in jene Regeln überzugehen und so Wiederholungen im Gesprächsablauf zu vermeiden \citep{chatscript2019}.\\


\section{Probleme mit der GUI}
\label{sec: Probleme mit der GUI}

Bei dem Testen der GUI fallen einige Probleme auf: 

\begin{enumerate}
\item{Fehlender Support von PostgreSQL-Server für Mac und Linux}
\item{Server mit PostgreSQL-Unterstüzung vs. Server ohne PostgreSQL-Unterstüzung}
\end{enumerate}

Das Testen der GUI erfolgte durch das Hosten auf Localhost (127.0.0.1) der durch ChatScript bereitgestellten Dateien im Ordner \textit{SERVER BATCH FILES} \citep{chatscript2019}. Während das Aufsetzen eines Servers auf Windows recht einfach war, fiel schnell auf, dass einige unvorhergesehene Probleme auftraten. 

\subsection{Fehlender Support von PostgreSQL-Server für Mac und Linux}
\label{sec:Fehlender Support von PostgreSQL-Server für Mac und Linux}

In der Dokumentation für ChatScript wird beschrieben, wie ein Server ohne Postgres-Unterstützung auf Windows, Mac und Linux-Distributionen aufzusetzen ist. Für Windows wird die .bat-Datei \textit{LocalPgServer.bat} mitgeliefert, sodass das Aufsetzen eines Servers auf Windows keine Probleme bereitet. Im Gegensatz dazu wird nichts über Server mit Postgres-Support für Linux oder Mac in der Dokumentation erwähnt. Nach mehreren erfolglosen Versuchen wurde im Forum für Chatbots \cite{cb5} nachgefragt, wobei diese Frage unbeantwortet blieb. 

\subsection{Server mit PostgreSQL-Unterstüzung vs. Server ohne PostgreSQL-Unterstüzung}
\label{sec: Server mit PostgreSQL-Unterstüzung vs. Server ohne PostgreSQL-Unterstüzung}

Ein weiteres Problem bei dem Testen der GUI war, dass der Bot verschiedene Verhaltensweisen je nach .bat-Datei aufweist. So entspricht der Programmablauf des Bots dem der Konsolenausgabe, wenn ein Server ohne Postgres-Support, zum Beispiel durch \textit{LocalServer.bat}, eingerichtet wird. Hier verhält sich der Bot in den Zustandsübergängen identisch zu der Konsolenausgabe, mit dem Unterschied, dass der Nutzer  bei der GUI / des Webinterfaces die erste Eingabe machen muss.\\
Das bedeutet, dass bis zum Zeitpunkt der Datenbankabfrage in Dbsearch keine Probleme auftreten. Zu Schwierigkeiten kommt es in diesem Fall verständlicherweise erst dann, wenn eine Datenbankabfrage ausgeführt werden muss, da hier keine Unterstützung für PostgreSQL besteht.\\
Im Gegensatz dazu gab es in vielerlei Hinsicht Probleme mit der Einrichtung eines Servers mit Postgres-Unterstüzung. Bei einigen Anläufen startete der Bot gar nicht, im besten Fall wird \texttildelow dbsearch.RESPONSE erreicht, jedoch geht der Bot von dort aus in keine weitere Regel, sondern iteriert weitergehend auf \texttildelow dbsearch.RESPONSE.\\
In vielen Versuchen erwiesen sich Übergänge in Regeln mit \lstinline|^reuse ()| als problematisch. Bei Aufsetzen des Servers mit \textit{LocalPgServer.bat}, ist es vorgekommen, dass einige der vorgesehenen Regelübergänge per \lstinline|^reuse()| nicht aktiviert wurden. \\
Da Probleme mit \textit{LocalPgServer.bat} nicht in der Dokumentation von ChatScript behandelt werden und der Macher und Entwickler von ChatScript, Bruce Wilcox, nicht auf Nachfragen mit dem Hosten eines Servers mit Postgres-Unterstützung antwortet \cite{cb5}, konnten alle Probleme mit dem Server mit Postgres-Unterstützung zeitnah nicht gelöst werden.\\
Es besteht unsererseits die Vermutung, dass es bei \textit{LocalPgServer.bat} Schwierigkeiten mit der User-Datenbank gibt, die ChatScript automatisch erstellt. In der Datenbank \textit{user} werden die einzelnen User mit Name und deren Protokolle abgespeichert. Diese Protokolle sind auch über den Ordner \textit{USER} zugänglich und tragen den Namen \textit{topic\_username\_botname.txt}. Interessanterweise wird bei \textit{LocalPgServer.bat} zwar ein Protokoll für den Benutzer angelegt, das existiert aber nur in der Datenbank und erscheint in dem USERS-Ordner nicht. Zudem bricht es nach SUMMARY ab, als wäre der Bot noch immer in diesem Zustand, obwohl die Datenbanksuche schon ausgeführt wurde und er aufgrund des ausgegebenen Textes schon in der Regel RESPONSE sein müsste.\\
Wie bereits erwähnt, tritt dieses Problem bei der Version über die Konsole und \textit{LocalSever.bat} nicht auf. Hier die Protokollführung nicht ab und die entsprechenden Protokolle sind über USERS einsehbar.\\
In der Datei \textit{bugs.txt} aus dem Ordner \textit{LOGS} wird das Problem folgendermaßen beschrieben: \lstinline{Postgres filessys write failed for USERS/topic\_127001\_augusta.txt}. Eine Suche über Google ergab, dass diese Beschreibung schon bei Schwierigkeiten mit ChatScript und dem Amazon Server aufgetaucht war (\url{https://www.chatbots.org/ai_zone/viewthread/3664/}). Die dortige Lösung das Problems, das Unique-Constraint der Relation \textit{userfiles} zu entfernen, haben wir auf unseren Fall übertragen und ausprobiert, allerdings hat dies nicht den gewünschten Effekt gehabt.\\
Aufgrund Zeitmangels, mangelnder Dokumentation und fehlender Einsicht in die Binärdateien bzw. Konfiguration ist dieses Problem bis dato nicht gelöst. Die genaue Fehlerquelle konnte nicht ausfindig gemacht werden und es ist ein Rätsel, warum auch Nicht-Postgres-Abläufe sich bei dem Server mit Postgres-Unterstützung anders verhalten als ohne Postgres-Unterstützung.

\chapter{Fazit}
\label{sec:Fazit}


%Glossar ausgeben
%\printglossary[style=altlist,title=Glossar]
\printglossary

\listoftables				% Tabellenverzeichnis
\listoffigures			% Abbildungsverzeichnis

\newpage

%\setbibpreamble{Präambel}		% Text vor dem Verzeichnis
%\bibliographystyle{dinat}
%\bibliography{bibliographie}	% Sie benötigen einen *.bib-Datei

%%%%%%%%%%%%%%%%%%%%%%%%%%%%%%%%%%%%%%%%%%%%%%%%
% B I B L I O G R A P H I E
%%%%%%%%%%%%%%%%%%%%%%%%%%%%%%%%%%%%%%%%%%%%%%%%

\renewcommand{\baselinestretch}{1}\normalsize 	%Zeilenabstand innerhalb der Angaben
\cleardoublepage
\addcontentsline{toc}{chapter}{\bibname}	%Bib. im Inhaltsvz.

\bibliography{CL-bib}					%Bibliographiedatenbank
\bibliographystyle{CL-DE}		%Zitierstruktur - natbib muss deklariert sein!

\nocite{*} 		%optional - gibt die gesammte Bib-Datenbank aus ansonsten werden nur 
%die zitierten Werke ausgegeben

%%%%%%%%%%%%%%%%%%%%%%%%%%%%%%%%%%%%%%%%%%%%%%%%
\end{document}
